%================================================================
\section{Relation to partial identification}\label{sec:pi}

The image borrows the reporting logic of partial identification, and it is worth
being precise about how far the analogy goes.

\paragraph{A menu of answers indexed by assumptions} In the partial-identification
tradition, an empirical question is answered by the set of values consistent with
the data and a stated list of assumptions. Weaker assumptions give larger sets and
more credible answers; stronger assumptions give smaller sets and less credible
ones, the trade that \citet{manski2003} calls the law of decreasing credibility.
\citet{tamer2010} describes the resulting practice as reporting a menu: one set
for each set of assumptions a reader might accept, instead of one point under the
strongest. This paper reports two such menus. The first is along the calibration
axis. The image $\Theta_{\hm}$ is the set of policy effects under the kernel, the
box, the targets, $\tau$ and $\delta$. Adding one more assumption, $c=c_0$,
collapses it to the single value $\theta_2(c_0)$; that is what the calibrator does
when she works at one point, and Table~\ref{tab:onepanel} shows the credibility
cost at the median-I point. The sensitivity menu of Section~\ref{sec:menu} moves
along the same axis by tightening $\delta$. The second menu is along the
interference axis. A0, A1 and A2 are three answers to the question of $\theta_2$
from one design, indexed by how much they assume about the way outcomes respond to
other people's seeding: no response at all, response through an exposure map with
transport, and the same with transport restricted to degree bins. Weaker
assumptions about interference buy less bias and cost precision, which is the same
trade in a different currency.

\paragraph{Where the set comes from} In much of the literature the set comes from
an incomplete model. A game with multiple equilibria does not say which outcome
occurs, so the model is consistent with a set of outcome distributions, and the
data bound the parameters without pinning them down \citep{tamer2003,tamer2010}.
Here the source is different. The kernel is complete: for each $c$ and each draw
of the noise it returns one panel, and the update map is a contraction with a
unique stationary regime (Appendix~\ref{app:model}). The set arises because the
targets do not restrict the parameters. Only $\alpha$ is pinned; the level
tolerance exceeds the pilot's own effect, so $\kappa$ stays open; and the
peer-correlation moment does not restrict $\gamma$. Many complete worlds fit the
same targets, and they disagree about the policy.

\paragraph{Where the analogy stops} The image is not an identified set. An
identified set is a population object: the values of a parameter consistent with
the distribution of the observables and the assumptions. The image here is a
finite empirical range produced by thresholding simulated moment discrepancies for
a fully specified kernel, evaluated on a sample of kept points. Its threshold mixes
Monte Carlo variance with a design tolerance, so it does not correspond to a
sampling distribution of the targets. A thresholded criterion resembles the level
sets of \citet{chernozhukov2007}, and the range of $\theta_2$ over such a set is a
projection of the kind studied by \citet{imbens2004} and \citet{kaido2019}; see
\citet{molinari2020} for a survey. Those papers give coverage guarantees under
their assumptions. Nothing here claims one. Similarly, no analyst on the
interference ladder reports an identified set for $\theta_2$; each reports a point
estimate with a standard error under its maintained assumptions. A bounds analysis
of $\theta_2$ from the pilot in the spirit of \citet{manski2013} would be a
different and more conservative exercise. What the paper borrows is the rule for
reporting: do not replace a set by a convenient point.
